% Continuation of Appendix A.
\subsection{Safe deletion and rechecking}

\begin{proof}[Proof of the deletion clauses of
\cref{thm:raw-frontier-closure-characterization}.]
For an active \(s\in L\), equation \eqref{eq:t3-state} is componentwise
equality of
\(K_{L^{-s}}=(F_{L^{-s}},C_{L^{-s}})\) and
\(K_L=(F_L,C_L)\).  The two redundancy equalities therefore imply dynamic
equivalence and finite-horizon value preservation by the forward projection
argument.  Under detectability, the converse characterization in A.2 makes
them necessary for dynamic innovation equivalence.  Compressed-state equality
also preserves the stationary value and complete action comparison under the
discounted contraction.
\end{proof}

\begin{proof}[Proof of \cref{cor:stepwise-safe-compression}.]
At each accepted deletion, apply the sufficient direction to the current
library, then compose the resulting compressed-state equalities along the
trace.  Every endpoint remains source-feasible and has the source's productive
values.  Each step removes an active strategy, so the procedure terminates.
If a complete final scan finds no certified deletion, the endpoint is
one-deletion irreducible.  The trace certifies feasibility and value
preservation, not minimum resource burden.
\end{proof}

\subsection{Operational/generative value and descendant bound}
\label{app:proofs-value}

\begin{proof}[Proof of \cref{thm:value-decomposition}.]
Substitute \(\Omega_n=U_n-P_n\) into the definition of
\(\Delta_n^{\gen}\).  The two passive-value differences combine into
\(\Delta_n^{\op}\), and the remaining full-value difference is
\(\mathcal I_n\).  This is the pointwise identity
\eqref{eq:value-decomposition} at fixed \((n,b,L,s)\).
\end{proof}

\begin{proof}[Proof of \cref{cor:frontier-silent-insertion}.]
Equal frontiers make the passive recursions identical.  If closure is also
unchanged, the two realizable compressed states coincide, so
\cref{thm:raw-to-compressed-projection} identifies their full values.
\end{proof}

\begin{proposition}[Research-premium monotonicity under project dominance]
\label{prop:premium-monotonicity}
Let \(L^-\) and \(L^+\) have the same frontier and
\(C_{L^-}\subseteq C_{L^+}\).  Suppose every old project remains feasible and
its complete action value---including cost, duration, operating rewards,
completion law, and continuation value---is weakly higher at \(L^+\).  Then
\[
  \Omega_n(b,L^-)\le\Omega_n(b,L^+)
\]
for every finite horizon and belief.  Hence a frontier-silent insertion
satisfying this project-dominance certificate has nonnegative generative
insertion value.
\end{proposition}

\begin{proof}
Induct on the finite horizon.  Continue has the same passive reward because
the frontiers agree.  Every old research action remains feasible and has a
weakly higher complete action value by hypothesis, while closure enrichment
may add actions.  Maximization therefore weakly raises full value.  Subtracting
the common passive value gives the premium order.
\end{proof}

\begin{proof}[Proof of \cref{thm:generative-carrier-lower-bound}.]
Condition on \((B_0,K_0,q)=(b,K_{L^+},q)\).  Full value at \(L^+\) dominates
both Continue and commitment to the enabled project \(q\).  Subtracting
passive value from the commitment value leaves initiation cost, the exact
operating adjustment \(A^{\op}_{q,h}\), and discounted continuation gain.
Insertion-only updates and retention of the passive policy make every omitted
outcome's continuation contribution nonnegative.  On the distinguished event
the gain is at least \(G(B_d)\), whose contribution under the joint completion
law is
\[
 \sum_{b'}\eta_{q,g}(b'\mid b,K_{L^+})G(b').
\]
Frontier silence removes operational insertion value, and the zero-premium
deleted comparator converts the retained premium into generative insertion
value.  Maximizing research against Continue yields
\eqref{eq:t6-lower-bound-adjusted}; setting \(A^{\op}=0\) gives
\eqref{eq:t6-lower-bound}.  Under the product specialization, substitution of
\(\eta_{q,g}=\pi\rho^d\mu_{q,d}\) gives
\(\pi\rho^d\overline G_{q,d}(b)\).  The sign comparisons in
\cref{cor:t6-comparative-statics} follow termwise from nonnegative joint mass
and gain, subject to its stated qualification when duration changes other
terms.
\end{proof}

\begin{example}[Exact one-belief retained carrier]
\label{ex:t6-exact-carrier}
Let \(d=1\), \(\beta=1/2\), \(\kappa=0\), and \(\pi=\rho=1\).  A zero-profile
bridge uniquely enables a descendant with gain \(G=2\); deletion leaves no
research project.  Operation continues against a zero frontier, so
\(A^{\op}=0\), and the bound gives
\[
  1\le\Delta^{\gen}_2.
\]
\end{example}

The extended gap-sum, coverage, and comparative-static arguments are in
Online Supplements~S1 and S3.
